\documentclass[../../../uem_phy_std_a.tex]{subfiles}

\begin{document}
    \begin{enumerate}

        \item まず，
            
            \myspaceb%
            $\displaystyle%
                x_1 = \uwave{Q_0} \;.
            $\\
            また，Yについてのキルヒホッフ則より，
            
            \myspaceb%
            $\displaystyle%
                -x_1 + y_1 = 0 \qquad \therefore \; y_1 = \uwave{Q_0} \;.
            $            
        \item Yについての電荷保存則より，
            
            \myspaceb%
            $\displaystyle%
                -x_2 + y_2 = 0 \;.
            $\\
            キルヒホッフ則より，
            
            \myspaceb%
            $\displaystyle%
                \myfracc{y_2}{\varepsilon_0 S} \cdot 2d + %
                \myfracc{x_2}{\varepsilon_0 S} \cdot d = V_0%
                \qquad \therefore \; x_2 = y_2%
                = \uwave{\myfracc{\varepsilon_0 S}{3d} V_0} \;.
            $
        \item Yについての電荷保存則より，
            
            \myspaceb%
            $\displaystyle%
                -x_3 + y_3 = Q_0 \;.
            $\\
            キルヒホッフ則より，
            
            \myspaceb%
            $\displaystyle%
                \myfracc{y_3}{\varepsilon_0 S} \cdot 2d + %
                \myfracc{x_3}{\varepsilon_0 S} \cdot d = 0
            $

            \myspaceb%
            $\displaystyle%
                \therefore \; x_3 = \uwave{- \myfraca{2}{3} Q_0} \;,%
                \quad y_3 = \uwave{\myfraca{1}{3} Q_0} \;.
            $
    \end{enumerate}
\end{document}
